\chapter{Допуск минимального вспомогательного фермионного модуля}
\label{ch:v9-kms-logdet-aux-fermion-admission}

\section{Наследуемые type- и package-носители}

Five-channel target-space уже имеет каноническое разложение
\begin{equation}
 V_{\rm type}=\mathbb C_s\oplus\mathbb C_a\oplus\mathbb C_t^3,
 \qquad \dim_{\mathbb C}V_{\rm type}=1+1+3=5.
 \label{eq:v9-kms-aux-fermion-type-carrier}
\end{equation}
Две независимо возникшие KMS-shapes образуют package-space
\begin{equation}
 P_{\rm KMS}=\mathbb C_\theta\oplus\mathbb C_\kappa,
 \qquad \dim_{\mathbb C}P_{\rm KMS}=2.
 \label{eq:v9-kms-aux-fermion-package-carrier}
\end{equation}
Поэтому минимальный ungraded carrier determinant-механизма не нужно
задавать покомпонентно: он функториально равен
\begin{equation}
 G_{\rm aux}^{\rm ungr}=V_{\rm type}\otimes P_{\rm KMS},
 \qquad \dim_{\mathbb C}G_{\rm aux}^{\rm ungr}=5\cdot2=10.
 \label{eq:v9-kms-aux-fermion-ungraded-carrier}
\end{equation}

\section{Нечётное поднятие}

Для fermionic Gaussian требуется parity shift
\begin{equation}
 G_{\rm aux}=\Pi G_{\rm aux}^{\rm ungr},
 \qquad
 \Gamma_{\rm aux}=-I_{10},
 \qquad \Gamma_{\rm aux}^2=I_{10}.
 \label{eq:v9-kms-aux-fermion-parity}
\end{equation}
Package-projectors имеют вид
\begin{equation}
 P_\theta=\begin{pmatrix}I_5&0\\0&0\end{pmatrix},
 \qquad
 P_\kappa=\begin{pmatrix}0&0\\0&I_5\end{pmatrix},
 \qquad
 P_\theta+P_\kappa=I_{10},\quad P_\theta P_\kappa=0.
 \label{eq:v9-kms-aux-fermion-package-projectors}
\end{equation}
Их ранги равны
\begin{equation}
 \rank P_\theta=\rank P_\kappa=5.
 \label{eq:v9-kms-aux-fermion-package-ranks}
\end{equation}

\section{Квадратичный оператор}

На нечётном носителе действует even operator
\begin{equation}
 D_{\rm aux}=R_\theta\oplus R_\kappa,
 \qquad [D_{\rm aux},\Gamma_{\rm aux}]=0.
 \label{eq:v9-kms-aux-fermion-operator}
\end{equation}
Его determinant в точности воспроизводит обе multiplicity patterns:
\begin{equation}
 \det D_{\rm aux}
 =r_{\theta s}r_{\theta a}r_{\theta t}^{3}
  r_{\kappa s}r_{\kappa a}r_{\kappa t}^{3},
 \qquad \deg\det D_{\rm aux}=10.
 \label{eq:v9-kms-aux-fermion-determinant}
\end{equation}
Поскольку triplet-block каждого $R$ скалярен, для family action
$U_t\in U(3)$ выполняется
\begin{equation}
 [D_{\rm aux},
  (1_s\oplus1_a\oplus U_t)\otimes I_2]=0.
 \label{eq:v9-kms-aux-fermion-family-covariance}
\end{equation}
Package-swap $W(\theta\leftrightarrow\kappa)$ действует ковариантно:
\begin{equation}
 W D_{\rm aux}(R_\theta,R_\kappa)W^{-1}
 =D_{\rm aux}(R_\kappa,R_\theta).
 \label{eq:v9-kms-aux-fermion-package-swap}
\end{equation}

\section{Berezin completion}

Complex Gaussian использует независимые odd variables
\begin{equation}
 \psi\in G_{\rm aux},
 \qquad \bar\psi\in G_{\rm aux}^{*},
 \qquad
 N_{\rm complex\ pairs}=10,\quad N_{\rm odd\ coordinates}=20.
 \label{eq:v9-kms-aux-fermion-complex-pairs}
\end{equation}
Эквивалентная antisymmetric pairing matrix равна
\begin{equation}
 A_{\rm B}=
 \begin{pmatrix}
 0&D_{\rm aux}\\
 -D_{\rm aux}^{T}&0
 \end{pmatrix},
 \qquad A_{\rm B}^{T}=-A_{\rm B}.
 \label{eq:v9-kms-aux-fermion-berezin-pairing}
\end{equation}
При positive shapes она невырождена:
\begin{equation}
 \rank A_{\rm B}=20,
 \qquad
 \det A_{\rm B}=(\det D_{\rm aux})^2.
 \label{eq:v9-kms-aux-fermion-pairing-rank}
\end{equation}
После выбора complex Berezin orientation это даёт
\begin{equation}
 \int D\bar\psi D\psi\,
 e^{-\bar\psi D_{\rm aux}\psi}
 =\det D_{\rm aux}
 =\det R_\theta\,\det R_\kappa.
 \label{eq:v9-kms-aux-fermion-integral}
\end{equation}

\section{Отделение от physical state-space}

Auxiliary carrier добавляется к пространство конфигураций как graded direct sum,
а не как новые QMS-states:
\begin{equation}
 \mathcal X_{\rm super}
 =\mathcal H_{\rm cr}^{\bar0}\oplus G_{\rm aux}^{\bar1},
 \qquad
 \dim_{\mathbb C}\mathcal H_{\rm cr}=6,\quad
 \dim_{\mathbb C}G_{\rm aux}=10.
 \label{eq:v9-kms-aux-fermion-superconfiguration}
\end{equation}
Quadratic operator block-diagonal,
\begin{equation}
 \mathcal D_{\rm super}=0_{\mathcal H_{\rm cr}}\oplus D_{\rm aux},
 \qquad
 \mathcal D_{\rm super}\iota_{\rm phys}=0,
 \qquad N_{\rm new\ physical\ states}=0.
 \label{eq:v9-kms-aux-fermion-decoupling}
\end{equation}
Следовательно, creation-QMS, его state dimension $6$ и Liouvillian
rank/nullity $35/1$ не меняются.

\section{Origin boundary}

Ungraded carrier выводится из двух уже зарегистрированных данных:
\begin{equation}
 N_{\rm functorial\ ungraded\ carrier\ origin}=2/2
 \quad\{V_{\rm type},P_{\rm KMS}\}.
 \label{eq:v9-kms-aux-fermion-ungraded-origin}
\end{equation}
Но ни four-slot parent, ни endpoint-QMS пока не содержат правила, которое
назначает этому tensor product нечётную статистику и Berezin orientation:
\begin{equation}
 N_{\rm odd\ statistics\ origin}=0/1,
 \qquad
 N_{\rm Berezin\ measure\ origin}=0/1.
 \label{eq:v9-kms-aux-fermion-statistics-origin}
\end{equation}
Это согласуется с общей BRST/BV границей: ghost-like fields могут входить
в graded generating functional, не увеличивая physical Hilbert space, но
их differential и gauge origin должны быть предъявлены отдельно.

\section{ProofDSL certificate и итог}

Точный LCF-слой проверяет type multiplicities, package-projectors, parity,
family covariance, determinant, minimal degree, package-swap, physical
decoupling и rank Berezin pairing:
\begin{equation}
 N_{\rm ProofDSL\ obligations}=10/10,
 \qquad
 N_{\rm registered\ gates}=51,
 \qquad
 N_{\rm registered\ obligations}=392.
 \label{eq:v9-kms-aux-fermion-proofdsl}
\end{equation}
Итоговый ledger равен
\begin{equation}
 \begin{aligned}
 N_{\rm auxiliary\ module\ architecture}&=10/10,\\
 N_{\rm functorial\ carrier\ origin}&=2/2,\\
 N_{\rm new\ physical\ states}&=0,\\
 N_{\rm odd\ statistics\ origin}&=0/1,\\
 N_{\rm Berezin\ measure\ origin}&=0/1,\\
 N_{\rm physical\ four\ slot\ parent}&=0/1.
 \end{aligned}
 \label{eq:v9-kms-aux-fermion-final-ledger}
\end{equation}

\begin{theorem}[Минимальный auxiliary module admitted]
Нечётный carrier
$G_{\rm aux}=\Pi(V_{\rm type}\otimes P_{\rm KMS})$ имеет минимальную
complex dimension $10$, сохраняет type/family/package symmetries и несёт
block-operator $R_\theta\oplus R_\kappa$. Его Berezin completion имеет
rank $20$ и не добавляет физических состояний creation-QMS.
\end{theorem}

\begin{nogocriterion}[Parity shift не выводится из tensor product]
Функториальное построение ungraded carrier не назначает ему Grassmann
statistics. Пока не построен BRST/BV differential, gauge redundancy или
иной parent-механизм нечётности, auxiliary module является условно
admitted, а логарифм определителя parent не считается физически выведенным.
\end{nogocriterion}

\begin{protocol}\begin{enumerate}
\item type carrier: \textbf{$1+1+3=5$};
\item package carrier: \textbf{$2$};
\item auxiliary odd dimension: \textbf{$10$};
\item Berezin odd coordinates: \textbf{$20$};
\item physical state increment: \textbf{$0$};
\item ProofDSL: \textbf{$10/10$, registry $51/392$};
\item statistics/measure origin: \textbf{$0/2$};
\item следующий гейт: \textbf{fermion-statistics parent-origin}.
\end{enumerate}\end{protocol}

Машинный сертификат сохранён в
\path{s2t_v9_endpoint_creation_kms_logdet_auxiliary_fermion_module_admission_gate_results.json}.

\begin{thebibliography}{9}
\bibitem{v9-kms-aux-berezin}
G.~Grensing, M.~Nitschmann,
\emph{Berezin integration over anticommuting variables and cyclic
cohomology}, arXiv:hep-th/0401231.
\bibitem{v9-kms-aux-brst}
M.~Henneaux, C.~Teitelboim,
\emph{Quantization of Gauge Systems},
Princeton University Press, 1992.
\end{thebibliography}